\section{Related Work}
\label{sec:related-work}

This section situates the present study relative to four bodies of prior
work: Hawkes processes in market microstructure, limit-order-book (LOB)
queue-position and queue-depletion modeling, machine-learning approaches
to LOB forecasting, and the small literature that compares parametric
point-process models against machine-learned forecasts directly. A fifth
subsection covers the statistical-testing methodology this paper relies
on. Every citation below was independently verified to exist (arXiv,
publisher, SSRN, or Semantic Scholar listing) before inclusion; where a
bibliographic detail such as venue or peer-review status could not be
confirmed, this is noted explicitly rather than presented as certain.

\subsection{Hawkes processes in market microstructure}
\label{sec:rw-hawkes}

Multivariate Hawkes processes are a standard tool for modeling
self-exciting order-flow clustering in financial markets;
\citet{bacry2015hawkes} survey the framework broadly, including
volatility estimation, market stability, contagion, and full
order-book dynamics, and this paper's 10-event marked multivariate
Hawkes process with exponential kernels and a branching-ratio
subcriticality constraint (Section~\ref{sec:methodology-hawkes}) is a
direct descendant of that framework. \citet{bowsher2007modelling}
introduces ``generalised Hawkes'' models for market events allowing
inhibitory effects and cross-day dependence, an early instance of the
same modeling stance taken here -- treating quote and trade events as a
multivariate point process with a full intensity vector, rather than a
single univariate arrival process. \citet{hewlett2006clustering} applies
a bivariate Hawkes process to buy/sell trade arrivals in FX to predict
imbalance, a narrower ancestor of the 10-event taxonomy of
Section~\ref{sec:data-taxonomy} (workshop presentation, not confirmed as
peer-reviewed). \citet{bormetti2015modelling} uses multivariate Hawkes
factor models to capture cojumps across correlated assets, a precedent
for using a shared or constrained Hawkes structure to transport
excitation information across instruments -- directly relevant to the
cross-asset transfer ladder of Section~\ref{sec:results-transfer}, which
transfers a compressed parametric object ($\boldsymbol\Gamma$) across
BTC, ETH, SOL, and DOGE rather than pooling raw training data.

Two papers bear directly on the near-critical branching-ratio values
observed in this study ($\rho(\boldsymbol\Gamma) \in [0.93, 0.98]$
throughout Table~\ref{tab:learning-curve}).
\citet{hardiman2013critical} finds that price-change Hawkes kernels on
S\&P E-mini data (1998--2011) are close to criticality (branching ratio
$\approx 1$) as a persistent stylized fact, consistent with this paper's
own measured range rather than suggesting a calibration artifact. Set
against that, \citet{filimonov2015apparent} warns that regime mixtures
and poor timestamp quality can produce \emph{spuriously} high estimated
branching ratios -- ``apparent criticality'' where the true value is
substantially lower -- a methodological caution directly relevant to
this paper's own tie-break causality sensitivity finding
(Section~\ref{sec:data-tiebreak}): the large swings in
$\boldsymbol\Gamma$'s Frobenius norm under alternative same-millisecond
tie-break conventions are exactly the kind of calibration fragility
\citeauthor{filimonov2015apparent} describe, which motivated reporting
tie-rule sensitivity explicitly rather than a single point estimate for
$\boldsymbol\Gamma$.

\subsection{Limit-order-book queue-position and queue-depletion modeling}
\label{sec:rw-lob}

\citet{cont2014price} shows that order-flow imbalance (OFI) at the best
quotes has an approximately linear relationship with short-horizon price
changes across a broad cross-section of NYSE stocks -- the direct
ancestor of this paper's $M_0$ baseline feature
(Section~\ref{sec:methodology-ladder}), and consistent with $M_0$ alone
already reaching PR-AUC $\approx 0.63$ on the fixed test week
(Table~\ref{tab:learning-curve}). \citet{gould2016queue} similarly finds
L1 bid/ask queue imbalance to be a strong, near-linear predictor of the
next tick move, particularly for large-tick instruments, supporting this
paper's inclusion of absolute queue sizes and rolling percentile rank in
$M_0$. \citet{huang2015simulating} models the LOB as a Markov
queuing system whose event intensities depend on current queue state --
the closest prior framework to this paper's Target B (pure depth
depletion), though a state-Markovian rather than history-Hawkes
formulation: where \citet{huang2015simulating} conditions intensities on
current depth, $M_2$ here instead conditions on recent \emph{event
history} via self- and cross-excitation, a complementary rather than
competing account of the same phenomenon.
\citet{moallemi2017queue} models the economic value of queue position via
a cancellation intensity increasing in queue size, an antecedent of this
paper's \texttt{*\_qty\_down} cancellation-volume event types
(Section~\ref{sec:data-taxonomy}) (working paper; no confirmed
peer-reviewed publication located).

\subsection{Machine learning for LOB forecasting}
\label{sec:rw-ml}

$M_1$'s design -- an unconstrained gradient-boosted-tree model on a
hand-engineered multi-scale event-count feature bank -- sits within an
established line of ML-for-LOB work. \citet{kercheval2015modelling} is an
early instance using multi-class SVMs on LOB-level attributes to predict
mid-price direction, which $M_1$ updates with a modern nonlinear tree
ensemble. \citet{ntakaris2018benchmark} establishes the widely used
FI-2010 benchmark (10 days, 5 Nasdaq Nordic stocks, $\sim$4M samples,
3-class up/stationary/down labels), against which this paper's own
labeling protocol departs in kind: Target A/B
(Section~\ref{sec:methodology-target}) predicts depletion and
adverse-move \emph{events} specifically, rather than the field's dominant
3-class mid-price-direction convention. \citet{zhang2019deeplob} (DeepLOB)
and \citet{wallbridge2020translob} (TransLOB, arXiv preprint, no
confirmed peer-reviewed venue located) represent the deep-learning and
attention-based ends of the same ``unconstrained representation
learning'' side of this paper's central contrast; this paper uses
LightGBM rather than a CNN/LSTM/Transformer as its free-model baseline, a
defensible and considerably cheaper stand-in at this paper's much finer
L1/event-level (rather than book-snapshot-image) granularity.
\citet{sirignano2019universal} shows a pooled, multi-stock LSTM
outperforms asset-specific models and generalizes to held-out stocks --
evidence for a ``universal,'' data-learnable price-formation mechanism
along a different transfer axis than this paper's own cross-asset
transfer ladder (Section~\ref{sec:results-transfer}): where
\citeauthor{sirignano2019universal} pool raw training data across
instruments into one free model, this paper instead transfers a
compressed parametric object ($\boldsymbol\Gamma$) and re-profiles only
the local baseline intensity and marks.

\subsection{Direct comparisons of parametric point-process and ML forecasts}
\label{sec:rw-comparison}

This is the category where the literature search found the least direct
overlap with the present paper's research question, and that gap is
reported here rather than papered over with a loose fit. No paper was
located that runs the same head-to-head this paper runs: a constrained,
MLE-fit multivariate Hawkes process against an unconstrained tree
ensemble, on identical LOB targets and an identical underlying feature
basis, under a formal walk-forward significance protocol. The closest
matches found are as follows. \citet{shi2022state} builds a \emph{hybrid}
neural-Hawkes model (continuous-time LSTMs modeling per-event-type
intensities plus an event-state interaction mechanism) and reports it
outperforms both pure-stochastic and pure-deep-learning baselines on LOB
event-stream prediction -- the closest published instance of a controlled
Hawkes-vs-ML comparison found, but a fusion study arguing that
\emph{combining} the two architectures wins, not a controlled test of
whether the parametric constraint itself generalizes better, which is
this paper's question. \citet{nitturanantha2026forecasting} reports that
a sum-of-exponentials Hawkes process gives the best forecast among
unspecified ``competing models'' for near-term order-flow-imbalance
forecasting on NSE tick data; whether its competitor set includes
tree-ensemble or deep-learning baselines could not be confirmed from the
publicly available abstract, so this paper does not claim it as a
confirmed ML comparison. \citet{mei2017neural} is the foundational
neural-Hawkes paper underlying the line of work
\citeauthor{shi2022state} build on (not finance-specific); it is relevant
background for why $M_2$ in this paper is deliberately kept fully
parametric (closed-form exponential kernels, explicit MLE, explicit
$\boldsymbol\Gamma$) rather than neuralized -- a neural Hawkes intensity
would no longer isolate the specific falsification-first question this
paper asks (Section~\ref{sec:intro-falsification}), namely whether the
parametric \emph{constraint itself} helps or hurts generalization.

In summary: the literature contains many papers on parametric Hawkes
processes for order flow, many papers on ML/deep learning for LOB
prediction, and a small number of hybrid fusion papers showing that
combining the two architectures outperforms either alone. What is
largely absent is a paper that holds the prediction target and the
free-model architecture fixed and asks, as this paper does, whether a
Hawkes-constrained feature representation regularizes better than letting
a tree see the same raw feature bank directly -- an ablation-style
comparison rather than an architecture-fusion or bake-off study. The
learning-curve crossover search of Section~\ref{sec:results-learning-curve}
appears to be a genuinely underexplored angle on the Hawkes-vs-ML
question relative to the literature surveyed here.

\subsection{Statistical methodology}
\label{sec:rw-methodology}

The rolling Giacomini--White protocol used throughout
(Section~\ref{sec:methodology-gw}) is a direct application of
\citet{giacomini2006tests}, which extends Diebold--Mariano/West-style
predictive-ability testing to settings where forecasting models are
estimated or re-estimated on a rolling basis and potentially
misspecified -- exactly the walk-forward, daily-re-estimation design this
paper's GW tests use (Sections~\ref{sec:results-gw}
and~\ref{sec:results-m3-gw}). The Romano--Wolf stepdown FWER control
across the 18-specification secondary grid
(Section~\ref{sec:results-secondary-grid}) follows
\citet{romano2005stepwise}, which frames stepdown multiple-testing
control explicitly around the data-snooping problem in empirical finance
(many strategies or specifications tested on the same data), together
with the companion methodological paper \citet{romano2005exact}
developing the underlying exact/approximate stepdown procedures; the two
are standard to cite jointly, as is done here.
